
\section{Related Work}
\label{sec:related_work}

\noindent\textbf{Unified understanding and generation models.}
Integrating visual comprehension and image synthesis within a single architecture has progressed rapidly. Early-fusion models such as Chameleon~\citep{team2024chameleon} process interleaved image-text sequences through a single transformer but achieve limited quality on both tasks. Decoupled visual encoding improved performance substantially: Janus~\citep{wu2025janus} and Janus-Pro~\citep{chen2025januspro} employ separate understanding and generation encoders while retaining weight sharing. Show-o~\citep{xie2024showo} and VILA-U~\citep{wu2024vilau} unify visual representations at the token level through omni-attention and unified tokenization, respectively. Autoregressive approaches such as Emu3~\citep{wang2024emu3} and VARGPT~\citep{xianwei2025vargpt} demonstrate that next-token prediction alone can match specialized systems, while TokenFlow~\citep{qu2025tokenflow} bridges both tasks with a dual-codebook tokenizer. On the generation side, BLIP3-o~\citep{chen2025blip3o} integrates diffusion-based synthesis, BAGEL~\citep{deng2025bagel} leverages flow-matching generation, and MetaMorph~\citep{tong2024metamorph} explores instruction tuning for joint capabilities. Despite rapid architectural progress, these models rely on supervised post-training with curated data, leaving open the question of whether continued self-supervised improvement is possible.

\noindent\textbf{Role-based self-play and self-improvement.}
Several methods decompose a unified model into interacting roles trained through cooperative or adversarial dynamics. UniCorn~\citep{han2026unicorn} identifies a ``conduction aphasia'' phenomenon where models understand visual content but fail to translate understanding into faithful generation, and addresses this by partitioning a model into Proposer, Solver, and Judge roles for cognitive pattern reconstruction. The dedicated Judge component, however, introduces a separate evaluation objective whose accuracy is bounded by the model's own capabilities. UniGame~\citep{su2025unigame} applies a lightweight perturber (less than 1\% additional parameters) at the shared token interface to expose fragile representations, primarily improving robustness and consistency, though generation gains remain modest compared to understanding improvements. SEER~\citep{tang2026endogenous} introduces a two-stage endogenous loop that first activates the model's latent evaluation ability through RL with verifiable rewards on a small proxy task, then uses this activated self-evaluator as a reward signal for generation. This bootstrapping is sample-efficient but requires a task-specific activation phase. EvoLMM~\citep{thawakar2025evolmm} is conceptually close to our approach in its proposer-solver structure and continuous self-consistency rewards, but addresses only visual understanding without extending to generation. A common thread is the introduction of auxiliary components or architectural dependencies that limit portability across model families.

\noindent\textbf{Reconstruction and internal-gap methods.}
A second line of work couples understanding and generation through reconstruction objectives or by exploiting the empirical observation that understanding consistently outperforms generation. UAE~\citep{yan2025understandinggeneration} formalizes this through an auto-encoder lens with three-stage training: cold-start reconstruction, Generation-for-Understanding (encoder learns captions that maximize decoder reconstruction), and Understanding-for-Generation (decoder refined on encoder-produced captions). RecA~\citep{xie2025reconstruction} takes a more efficient approach, conditioning generation on the understanding encoder's internal embeddings and optimizing image reconstruction in approximately 27 GPU-hours, though this requires access to architecture-specific hidden representations. The Internal Gap method~\citep{han2025internalgap} directly leverages the understanding-generation asymmetry by using the stronger understanding module to score self-generated images for DPO-based post-training, observing a co-improvement effect where generation gains also benefit understanding. UniMRG~\citep{su2026unimrg} takes a complementary direction, enhancing understanding by training the model to produce multiple intrinsic representations (pixel, depth, segmentation) as auxiliary generation tasks.

\noindent\textbf{Reinforcement-based optimization.}
RL formulations cast post-training as policy optimization over both tasks. UniRL~\citep{mao2025unirl} performs iterative self-training with GRPO where the model generates images, answers diagnostic questions on them, and computes rewards against ground-truth, though the reward computation still requires annotations. CoRL~\citep{jiang2025coreinforcement} proposes two-stage RL with cycle-consistency and text-image matching rewards, yielding strong reasoning gains but adding training complexity. SUDER~\citep{hong2025dualselfrewards} avoids external rewards by exploiting the duality between understanding and generation, where each direction provides reward signals for the other via SimPO optimization. SRUM~\citep{jin2025srum} introduces fine-grained spatial reward decomposition with global (compositional correctness) and local (attribute fidelity) rewards forming spatially-aware reward maps. SILMM~\citep{qu2024silmm} provides a self-improvement baseline for compositional generation through a multi-stage pipeline of VQA-based self-feedback and DPO alignment, while ILLUME~\citep{wang2024illume} contributes a self-enhancing scheme where the model self-assesses text-image consistency of its own generations. VARGPT-v1.1~\citep{zhuang2025vargpt} and Skywork UniPic 2.0~\citep{wei2025skyworkunipic2} apply iterative RL but rely on external preference data (Midjourney, FLUX~\citep{flux2024}) or external reward evaluators (GPT-4), departing from the self-supervised setting.

\noindent\textbf{Our approach.}
Our framework differs from prior work along several axes. Unlike role-based methods that introduce auxiliary judges~\citep{han2026unicorn}, architecture-specific perturbers~\citep{su2025unigame}, or reward activation phases~\citep{tang2026endogenous}, we require no additional components beyond three lightweight LoRA adapters on a frozen backbone. Unlike reconstruction approaches~\citep{yan2025understandinggeneration, xie2025reconstruction} that depend on encoder-decoder symmetry or access to internal representations, we operate solely through text generation and image synthesis interfaces. Unlike reinforcement-based methods that require ground-truth annotations~\citep{mao2025unirl}, task-specific reward design~\citep{jiang2025coreinforcement}, external preference data~\citep{zhuang2025vargpt}, or external reward evaluators~\citep{wei2025skyworkunipic2}, we derive all training signals from self-consistency and token-level entropy over unlabeled images. The multi-scale generation assessment provides richer signal than single-scalar alternatives~\citep{hong2025dualselfrewards} and more principled coverage than manual granularity design~\citep{jin2025srum}. Most importantly, we validate one matched training framework across three architecturally distinct model families, providing strong evidence for model-agnostic self-supervised post-training.
